\documentclass[11pt]{article}
\usepackage[margin=1in]{geometry}
\usepackage{amsmath,amssymb,amsthm}
\newtheorem{theorem}{Theorem}
\newtheorem{lemma}{Lemma}
\newtheorem{remark}{Remark}
\title{An Exactly Solvable Endpoint Transition Model}
\author{}\date{October 5, 2026 --- paper proof}
\begin{document}\maketitle

For $v\in L^2(0,\infty)$ set $q(s)=\int_0^s v(u)\,du$, $b(s)=s^2/2-s^3/6$, and define the extended-real functional
\[
 \mathscr E(v)=\frac12\int_0^\infty v^2\,ds+
 \int_0^\infty\left(\sup_{s>0}\frac{q(s)+b(s)-X}{s}\right)_+dX.
\]
The supremum is measurable: for every fixed $X$, its expression is continuous in $s>0$, and one may take the supremum over positive rational $s$. The second integral is nonnegative and is allowed to be infinite.

\begin{theorem}[Exact optimum, uniqueness, and stability]
The unique minimizer is
\[
 v_*(s)=\begin{cases}
 -3/16,&0<s<3/2,\\
 s^2/4-s/2,&3/2\le s\le2,\\
 0,&s>2.
 \end{cases}
\]
Its minimum is $\mathscr E(v_*)=229/3840$. More strongly, every $v\in L^2(0,\infty)$ satisfies
\[
 \boxed{\mathscr E(v)-\frac{229}{3840}\ge\frac12\|v-v_*\|_{L^2}^2.}
\]
For the unperturbed choice $v=0$ the value is $229/1920$, exactly twice the minimum.
\end{theorem}
\begin{proof}
Integration gives
\[
 q_*(s)=\begin{cases}
 -3s/16,&0\le s\le3/2,\\
 s^3/12-s^2/4,&3/2\le s\le2,\\
 -1/3,&s\ge2.
 \end{cases}
\]
Put $V=q_*+b$. On $[3/2,2]$ one has
\[
 V=s^2/4-s^3/12,\quad V'=s/2-s^2/4,\quad
 X(s)=V-sV'=s^2(2s-3)/12,\quad X'(s)=s(s-1)/2.
\]
Thus $X$ maps $[3/2,2]$ increasingly onto $[0,1/3]$.

For a fixed $X\ge0$, the derivative of $(V(s)-X)/s$ is $(X-[V-sV'])/s^2$. On $(0,3/2)$ the bracket is $-s^2/2+s^3/3\le0$, so this ratio increases there. On $[3/2,2]$ it reaches its maximum at the unique $s$ with $X=X(s)$ when $0\le X\le1/3$. On $(2,\infty)$, $V-sV'=-1/3-s^2/2+s^3/3$ increases from $1/3$, so no larger value occurs for those $X$. For $X>1/3$, $V(s)\le1/3$ for all $s$, hence the positive part of the supremum is zero. The maximizing envelope is consequently $V'(s(X))$ on $[0,1/3]$ and zero afterwards.

For an arbitrary competitor write $Q=q-q_*$. Evaluating its supremum at these same maximizing parameters yields the global affine lower bound
\begin{align*}
 &\int_0^\infty\left(\sup_{s>0}\frac{q(s)+b(s)-X}{s}\right)_+dX
 -\int_0^\infty\left(\sup_{s>0}\frac{q_*(s)+b(s)-X}{s}\right)_+dX\\
 &\hspace{12mm}\ge\int_0^{1/3}\frac{Q(s(X))}{s(X)}\,dX
 =\int_{3/2}^2\frac{s-1}{2}Q(s)\,ds.
\end{align*}
This remains valid if the competitor integral is infinite. Since $v_*$ is continuous, vanishes after two, and has derivative $(s-1)/2$ on $(3/2,2)$ and zero elsewhere, integration by parts gives
\[
 \int_0^\infty v_*(v-v_*)\,ds=-\int_{3/2}^2\frac{s-1}{2}Q(s)\,ds.
\]
There is no boundary term at zero because $Q(0)=0$, or at infinity because $v_*$ vanishes there. Expanding the kinetic square cancels the affine lower bound and proves the stated exact rigidity inequality.

Let $h(s)=s/2-s^2/4$. The kinetic part at the minimizer and its envelope part both equal
\[
 \frac12\left(\frac32(3/16)^2+\int_{3/2}^2h(s)^2\,ds\right)=\frac{229}{7680}.
\]
For the envelope part, use $X'=-s h'$ and integration by parts. These two values add to $229/3840$. When $v=0$, the exposed interval is again $[3/2,2]$, but the exposed height is $s-s^2/2=2h(s)$. Its area is four times the minimizer's envelope area, and there is no kinetic term. Thus $\mathscr E(0)=229/1920$.
\end{proof}

\section*{Why this model occurs, and the transfer obligation}
Write $\omega=1+\delta$, and let $p_\delta^*$ denote Baek's relaxed support. Perturb it near the right endpoint by
\[
 p(t)=p_\delta^*(t)+\delta^3P(t/\delta),\qquad q(s)=P(s)-P(0).
\]
The floor intercept relative to its initial value has the expansion
\[
 \frac{p(\delta s)-1}{\cos(\delta s)}-(p(0)-1)
 =\delta^3\left(q(s)+\frac{s^2}{2}-\frac{s^3}{6}\right)+o(\delta^3)
\]
on compact $s$ intervals, for fixed sufficiently regular $P$. The tent height at abscissa $p(0)-1+\delta^3X$ consequently scales like $\delta^2(q(s)+b(s)-X)/s$. Its area scale is $\delta^5$. The support-gap energy has the same scale, with leading term $\tfrac12\delta^5\int P'^2$.

The two endpoints suggest twice the model minimum, namely $229/1920$, as a candidate coefficient of the unrestricted deficit below $1+\omega^2/2$. This calculation alone is not a theorem about that deficit. Such a theorem needs a compactness/lower-bound argument for arbitrary maximizing caps, separation of the two endpoint contributions, an admissible support recovery sequence, and control of any remaining niche excess over the core.

There is also a real endpoint issue: $v_*(0+)=-3/16$, whereas a smooth normalized cap without an endpoint atom has $p'(0)=0$. The model must therefore be recovered with a thinner boundary layer; directly inserting its derivative into a cap violates the endpoint condition. This obstacle is recorded rather than silently ignored.

\paragraph{Validation.}
The proof is an exact subgradient certificate. The polynomial integrals were also checked symbolically, but their displayed antiderivatives and the certificate suffice. No CI or compilation was run. Until the transfer obligations are supplied separately, this note claims only the model theorem, not closure of the larger-angle sofa problem.
\end{document}
